\part{Worked Examples and the Admissibility Layer}

\secn{10}{10}{Admissible exact readings}

The examples specialize the exact kernel to diagnosis, database dependencies, and finite-machine observability. The admissibility layer then separates raw content, joint admissible content, and the content attainable by separate per-view readings. Monotone ceilings guarantee a bracket; separability is the stronger join-preservation property. Adjoint claims use the full subset lattice, with the finite-view convention or the explicit arbitrary-join extension in \hyperlink{stmt:15.4}{Corollary~15.4}.

\secn{11}{11}{Worked example A: clinical diagnosis and missing information}

This example exhibits, on the smallest scale that shows all phenomena, the full cycle: compute the induced quotient, detect an unanswerable question via unresolved pairs, and verify the marginal value of a new view.

\textbf{Setup.} Let the contrast domain be six patient states,
\[
D=\{s_1,s_2,s_3,s_4,s_5,s_6\}.
\]
Three views are available: a lab reading \(L:D\to\{\mathrm{lo},\mathrm{hi}\}\), a symptom flag \(F:D\to\{0,1\}\), and an age bracket \(A:D\to\{\mathrm{y},\mathrm{o}\}\). The question is a diagnosis \(q:D\to\{\text{well},\text{mild},\text{severe}\}\). The data are given by the table:

{\def\LTcaptype{none} % do not increment counter
\begin{longtable}[]{@{}
  >{\raggedright\arraybackslash}p{(\linewidth - 8\tabcolsep) * \real{0.2000}}
  >{\raggedright\arraybackslash}p{(\linewidth - 8\tabcolsep) * \real{0.2000}}
  >{\raggedright\arraybackslash}p{(\linewidth - 8\tabcolsep) * \real{0.2000}}
  >{\raggedright\arraybackslash}p{(\linewidth - 8\tabcolsep) * \real{0.2000}}
  >{\raggedright\arraybackslash}p{(\linewidth - 8\tabcolsep) * \real{0.2000}}@{}}
\toprule\noalign{}
\begin{minipage}[b]{\linewidth}\raggedright
state
\end{minipage} & \begin{minipage}[b]{\linewidth}\raggedright
\(L\)
\end{minipage} & \begin{minipage}[b]{\linewidth}\raggedright
\(F\)
\end{minipage} & \begin{minipage}[b]{\linewidth}\raggedright
\(A\)
\end{minipage} & \begin{minipage}[b]{\linewidth}\raggedright
\(q\)
\end{minipage} \\
\midrule\noalign{}
\endhead
\bottomrule\noalign{}
\endlastfoot
\(s_1\) & lo & 0 & y & well \\
\(s_2\) & lo & 0 & o & mild \\
\(s_3\) & lo & 1 & y & mild \\
\(s_4\) & hi & 0 & y & mild \\
\(s_5\) & hi & 1 & y & severe \\
\(s_6\) & hi & 1 & o & severe \\
\end{longtable}
}

\textbf{Content of \(\{L,F\}\).} The view partitions are
\[
\ker(L)=\{\,s_1s_2s_3\mid s_4s_5s_6\,\},\qquad
\ker(F)=\{\,s_1s_2s_4\mid s_3s_5s_6\,\}.
\]
Their join (common refinement, \hyperlink{stmt:3.3}{Fact~3.3}), which by \hyperlink{stmt:5.2}{Definition~5.2} is \(\sigma(\{L,F\})\), groups states by the pair \((L,F)\):
\[
\sigma(\{L,F\})=\{\,s_1s_2\mid s_3\mid s_4\mid s_5s_6\,\}.
\]

\textbf{Answerability of \(q\).} By \hyperlink{stmt:5.6}{Proposition~5.6}, \(q\) is answerable from \(\{L,F\}\) iff it is constant on every block. It is constant on \(\{s_3\}\), \(\{s_4\}\), and \(\{s_5s_6\}\) (both severe), but on the block \(\{s_1s_2\}\) we have \(q(s_1)=\text{well}\ne\text{mild}=q(s_2)\). Hence \(q\) is \textbf{not} answerable from \(\{L,F\}\), and the unresolved-pair set (Definition 5.5) is
\[
U(\{L,F\},q)=\{(s_1,s_2),(s_2,s_1)\}.
\]
The states \(s_1,s_2\) agree on both available views yet differ in diagnosis: this is the exact locus of missing information.

\textbf{Marginal value of \(A\).} The remedy must separate the unresolved pair. Since \(A(s_1)=\mathrm{y}\ne\mathrm{o}=A(s_2)\), the view \(A\) separates \((s_1,s_2)\), so by \hyperlink{stmt:6.5}{Theorem~6.5}(3) adding \(A\) makes \(q\) answerable. Indeed the profiles \((L,F,A)\) are pairwise distinct across all six states, so
\[
\sigma(\{L,F,A\})=\top\quad(\text{the discrete partition}),
\]
and \(q\) factors trivially. Note the audit did not require examining \(q\) against every subset of views: by \hyperlink{stmt:6.8}{Corollary~6.8} it reduced to comparing one partition, \(\sigma(S)\), against \(\ker(q)\), and the fix was dictated by one pair.

The example also illustrates why content is domain-relative (Definition 4.1): had \(D\) excluded \(s_2\), the pair would not exist and \(\{L,F\}\) would already answer \(q\). Sufficiency is a property of data \emph{against a contrast domain and a question}, never of data alone.

\secn{12}{12}{Worked example B: relational dependency theory recovered}

This example validates the kernel against an established theory: functional dependencies in the relational model. The claim is that answerability, the induced quotient, and the universal property specialize exactly to functional determination, the agreement partition, and projection---\hspace{0pt}and that Armstrong's inference rules become lattice identities.

\textbf{Setup.} Fix a relation instance: a finite set \(D\) of tuples over an attribute set \(\mathcal U\). For \(X\subseteq\mathcal U\), the \textbf{projection view} is \(\pi_X:D\to\prod_{a\in X}\mathrm{dom}(a)\), \(t\mapsto t{\restriction}X\). Its kernel \(\ker(\pi_X)\) is the classical \emph{agreement partition}: two tuples share a block iff they agree on every attribute in \(X\).

\textbf{Determination is answerability.} A functional dependency \(X\to Y\) holds in the instance iff any two tuples agreeing on \(X\) agree on \(Y\), i.e.~iff \(\pi_Y\) is constant on the blocks of \(\ker(\pi_X)\). By \hyperlink{stmt:3.2}{Fact~3.2} and \hyperlink{stmt:5.4}{Definition~5.4} this is exactly
\[
X\to Y \quad\Longleftrightarrow\quad \pi_Y \text{ answerable from } \{\pi_X\}
\quad\Longleftrightarrow\quad \ker(\pi_Y)\le\ker(\pi_X),
\]
that is, the \(X\)-partition refines the \(Y\)-partition. Determination \emph{is} answerability; the ``closer to determined'' ordering is refinement.

\textbf{Content of a projection set.} For a set of attributes viewed as \(S=\{\pi_{X_1},\dots,\pi_{X_k}\}\), \hyperlink{stmt:5.2}{Definition~5.2} and (3.1) give
\[
\sigma(S)=\bigvee_i \ker(\pi_{X_i})=\ker(\pi_{X_1\cup\cdots\cup X_k}),
\]
the agreement partition of the union of attributes. By \hyperlink{stmt:6.7}{Theorem~6.7}(3), the induced quotient is (isomorphic to) the image of the joint projection \(\pi_{\bigcup_i X_i}\): the set of distinct \((\bigcup_i X_i)\)-value profiles realized in the instance. The kernel's ``free lossless summary of the data'' is, here, the projected relation.

\textbf{Armstrong's axioms as lattice facts.} Writing \(\ker_X:=\ker(\pi_X)\) and using \(\ker_{X\cup Z}=\ker_X\vee\ker_Z\):

\begin{thmplain}[\hypertarget{stmt:12.1}{Proposition~12.1}\ {\normalfont\upshape(soundness of the Armstrong rules)}]

The three Armstrong rules hold as identities in \(\mathrm{Part}(D)\):

\begin{itemize}
\tightlist
\item
  \emph{Reflexivity:} if \(Y\subseteq X\) then \(\ker_Y\le\ker_X\), so \(X\to Y\).
\item
  \emph{Augmentation:} if \(\ker_Y\le\ker_X\) then \(\ker_{Y\cup Z}=\ker_Y\vee\ker_Z\le\ker_X\vee\ker_Z=\ker_{X\cup Z}\), so \(XZ\to YZ\).
\item
  \emph{Transitivity:} if \(\ker_Z\le\ker_Y\) and \(\ker_Y\le\ker_X\) then \(\ker_Z\le\ker_X\), so \(X\to Z\).
\end{itemize}
\end{thmplain}

\begin{proof}

Reflexivity: \(\pi_X\) pairs \(\pi_Y\) with \(\pi_{X\setminus Y}\), so \(\ker_X=\ker_Y\vee\ker_{X\setminus Y}\ge\ker_Y\). Augmentation and transitivity are monotonicity of \(\vee\) and transitivity of \(\le\), as displayed.
\end{proof}

\textbf{Mathematical scope.} These are sound rules for the fixed view family. Completeness of Armstrong inference over all relational models is a separate classical theorem; the displayed lattice proof does not prove that every dependency true in this particular family is derivable from an arbitrarily chosen dependency basis.

Thus the completeness of Armstrong's system, in a fixed instance, reflects that \(\{\ker_X\}_{X\subseteq\mathcal U}\) is a join-subsemilattice of \(\mathrm{Part}(D)\) and that functional determination is its order. The partition view of dependencies is the basis of dependency-discovery algorithms \autocite{huhtala1999}, and its appearance here as a special case is evidence that the kernel's primitives are the right ones: a theory built for heterogeneous views specializes, with no adjustment, to the attribute-incidence case that classical dependency theory and Formal Concept Analysis already treat \autocite{armstrong1974,ganter1999}.

\secn{13}{13}{Worked example C: observability of a finite machine}

This example validates the kernel against control- and automata-theoretic observability (cf.~Kalman 1963; Nerode 1958): the induced quotient of the output-history views is the observable-state partition, and answerability of ``which state?'' is observability.

\textbf{Setup.} Let a deterministic Moore machine have state set \(D=\{p,q,r\}\), a single input letter with transition \(\delta:D\to D\) given by \(p\mapsto q\), \(q\mapsto r\), \(r\mapsto r\), and output \(\lambda:D\to\{0,1\}\) with \(\lambda(p)=\lambda(q)=0\), \(\lambda(r)=1\). The \textbf{depth-\(k\) view} is the output observed after \(k\) steps,
\[
v_k:=\lambda\circ\delta^{\,k}:D\to\{0,1\},\qquad k=0,1,2,\dots
\]
Each \(v_k\) is a legitimate view on \(D\); the available family after observing horizon \(H\) is \(S_H=\{v_0,\dots,v_H\}\).

\textbf{Computing the observable quotient.}
\[
v_0:\ p\mapsto0,\ q\mapsto0,\ r\mapsto1
\ \Rightarrow\ \ker(v_0)=\{\,pq\mid r\,\};
\]
\[
v_1=\lambda\circ\delta:\ p\mapsto\lambda(q)=0,\ q\mapsto\lambda(r)=1,\ r\mapsto\lambda(r)=1
\ \Rightarrow\ \ker(v_1)=\{\,p\mid qr\,\}.
\]
Then
\[
\sigma(S_1)=\ker(v_0)\vee\ker(v_1)=\{\,pq\mid r\,\}\vee\{\,p\mid qr\,\}=\{\,p\mid q\mid r\,\}=\top .
\]
The three states become pairwise distinguishable at horizon \(1\): the machine is observable, and the question ``which state?'' (the identity question \(\mathrm{id}:D\to D\), with \(\ker(\mathrm{id})=\top\)) is answerable from \(S_1\) by \hyperlink{stmt:5.6}{Proposition~5.6}.

\textbf{Reading the kernel notions in automata terms.} The induced quotient \(\sigma(S_H)\) is exactly the partition that identifies states not separable within horizon \(H\); its stabilization \(\sigma(S_\infty)\) is the Nerode/observability partition, and the machine is minimal/observable iff \(\sigma(S_\infty)=\top\). Unresolved pairs (Definition 5.5) for the identity question are precisely the \emph{indistinguishable state pairs}; \hyperlink{stmt:6.5}{Theorem~6.5} says a longer horizon helps exactly when a newly observed output separates some still-merged pair---\hspace{0pt}the partition-refinement step of state-minimization algorithms \autocite{hopcroft1979}. The observability special case of the kernel thus coincides with the classical construction, and control-theoretic observability \autocite{kalman1963} is its dynamical-state instance.

Taken together, Examples A--C show the kernel is neither vacuous nor ad hoc: two of the three instances recover established theories on the nose, and the third yields the operationally correct diagnosis-and-remedy. With the kernel thus anchored, the remainder of Part II relaxes canonical registration.

\secn{14}{14}{The admissibility layer: registered partitions and admissible ceilings}

Standing assumption (A3) fixed registration to the canonical \(\kappa^{\mathrm{can}}_v(y)=v^{-1}(y)\), which by \hyperlink{stmt:4.7}{Lemma~4.7} is the most informative sound registration. Real interpreters---\hspace{0pt}bounded decoders, invariant readouts, resource-limited extractors---\hspace{0pt}cannot always realize it. This section relaxes (A3): registration is henceforth chosen from an admissibility class, and the content it induces is studied as that class varies. The development stays exact and deterministic; assumptions (A1), (A2), (A4), (A5) remain in force.

\begin{quote}
\textbf{Convention (finiteness of view sets).} From this point through the end of Part III, all view sets \(S, T \subseteq V\) are finite unless explicitly stated otherwise: \hyperlink{stmt:14.3}{Definition~14.3} assigns ceilings to finite sets only, and the joint notions below are defined through them. The one systematic exception is recorded where it occurs---\hspace{0pt}the separable adjunction (Theorem 15.1) involves only singleton ceilings and holds for arbitrary \(S \subseteq V\).
\end{quote}

\subsecn{1}{14.1}{The partition induced by a registration}

\begin{thmdef}[\hypertarget{stmt:14.1}{Definition~14.1}]

For a view \(v\) and a registration \(\kappa_v:Y_v\to\Phi_D\), the \textbf{registered partition} is
\[
\rho(\kappa_v):=\ker(\kappa_v\circ v)\in\mathrm{Part}(D),
\]
the partition of \(D\) by equality of registered pieces: \(Z\sim Z'\) iff \(\kappa_v(v(Z))=\kappa_v(v(Z'))\).
\end{thmdef}

\begin{thmplain}[\hypertarget{stmt:14.2}{Lemma~14.2}\ {\normalfont\upshape(registration is distinction--non-increasing)}]

For every registration \(\kappa_v\),
\[
\rho(\kappa_v)\le\ker(v),
\]
with equality for the canonical registration: \(\rho(\kappa^{\mathrm{can}}_v)=\ker(v)\).
\end{thmplain}

\begin{proof}

The map \(\kappa_v\circ v\) factors through \(v\) by construction, so \(v(Z)=v(Z')\Rightarrow\kappa_v(v(Z))=\kappa_v(v(Z'))\); hence \(\ker(v)\) refines nothing beyond \(\rho(\kappa_v)\), i.e.~\(\rho(\kappa_v)\le\ker(v)\). For canonical, \(\kappa^{\mathrm{can}}_v(v(Z))=v^{-1}(v(Z))=[Z]_{\ker(v)}\), so the induced partition is \(\ker(v)\) itself.
\end{proof}

\hyperlink{stmt:14.2}{Lemma~14.2} is the exact counterpart, at the level of distinctions, of \hyperlink{stmt:4.7}{Lemma~4.7}'s statement at the level of pieces: registration can only \emph{discard} the distinctions a view already draws, never manufacture new ones, and canonical registration discards none. Everything downstream lives in the interval \([\bot,\ker(v)]\).

\subsecn{2}{14.2}{Admissibility structures}

The kernel used one registration per view; admissibility allows a \emph{class} and, crucially, must also say what is admissible on \emph{compound} views, because that is where the theory's new content lies. We formulate admissibility directly at the level of registered partitions, since by \hyperlink{stmt:14.2}{Lemma~14.2} the partition is what determines content.

\begin{thmdef}[\hypertarget{stmt:14.3}{Definition~14.3}]

For a view family \(V\) on \(D\), an \textbf{admissibility structure} \(K\) assigns to each finite \(T\subseteq V\) an \textbf{admissible ceiling}
\[
\gamma_K(T)\in\mathrm{Part}(D),
\]
subject to:

\begin{itemize}
\tightlist
\item
  \textbf{(K0) Grounding.} \(\gamma_K(\varnothing)=\bot\).
\item
  \textbf{(K1) Boundedness.} \(\gamma_K(T)\le\sigma(T)\) for all \(T\) (no registration invents distinctions its raw material does not draw; \hyperlink{stmt:14.2}{Lemma~14.2}).
\item
  \textbf{(K2) Realizability closure.} For each \(T\), every \(\rho\le\gamma_K(T)\) is itself admissible on \(T\); i.e.~an admissible reading may always be coarsened. Equivalently, the admissible partitions on \(T\) form the principal ideal \({\downarrow}\gamma_K(T)\).
\end{itemize}

\(K\) is \textbf{monotone} if \(T\subseteq T'\Rightarrow\gamma_K(T)\le\gamma_K(T')\): enlarging the raw material never lowers the admissible ceiling.
\end{thmdef}

(K2) records that admissibility caps \emph{resolving power}, not the freedom to ignore: given an admissible reading one may always throw information away, so the admissible partitions on a fixed \(T\) are exactly those below a top element \(\gamma_K(T)\). The content of an admissibility structure therefore lies not in the per-\(T\) ideals but in \emph{how the ceilings on different \(T\) relate}---\hspace{0pt}the assignment \(T\mapsto\gamma_K(T)\) as a whole. Monotonicity is optional and marks a fault line examined in \hyperlink{sec:16}{§16}.

That the axioms describe realizable structure, rather than stipulating it, is guaranteed by the following lemma.

\begin{thmplain}[\hypertarget{stmt:14.4}{Lemma~14.4}\ {\normalfont\upshape(realization of admissible partitions)}]

Let \(T \subseteq V\) be finite, let \(c_T = \langle v \rangle_{v \in T} : D \to \prod_{v \in T} Y_v\) be the compound view, and let \(\rho \le \sigma(T)\). Then the registration defined on realized profiles by
\[
\kappa\big(c_T(Z)\big) = [Z]_\rho
\]
(extended by the null value \(D\) off the image) is well defined, sound, and has registered partition \(\rho(\kappa) = \rho\). Together with \hyperlink{stmt:14.2}{Lemma~14.2} applied to the compound view, the partitions realizable by sound registrations of \(c_T\) are therefore exactly the principal ideal \({\downarrow}\,\sigma(T)\): axioms (K1)--(K2) delimit precisely the realizable range, and nothing in \hyperlink{stmt:14.3}{Definition~14.3} is stipulative beyond the choice of ceiling.
\end{thmplain}

\begin{proof}

Well-definedness: \(c_T(Z) = c_T(Z')\) means \(Z \sim_T Z'\), i.e.~\([Z]_{\sigma(T)} = [Z']_{\sigma(T)}\); since \(\rho \le \sigma(T)\), every \(\sigma(T)\)-block lies inside a \(\rho\)-block, so \([Z]_\rho = [Z']_\rho\). Soundness: \(Z \in [Z]_\rho\). Registered partition: \(\kappa(c_T(Z)) = \kappa(c_T(Z'))\) iff \([Z]_\rho = [Z']_\rho\), so \(\rho(\kappa) = \ker(Z \mapsto [Z]_\rho) = \rho\). The converse bound \(\rho(\kappa) \le \ker(c_T) = \sigma(T)\) for every registration is \hyperlink{stmt:14.2}{Lemma~14.2}.
\end{proof}

\hyperlink{stmt:14.4}{Lemma~14.4} grounds each ceiling \emph{individually}. Whether a family of ceilings is \emph{simultaneously} realizable by one coherent class of decoders is a further question, and it is exactly where monotonicity acquires operational content.

\begin{thmplain}[\hypertarget{stmt:14.4′}{Proposition~14.4′}\ {\normalfont\upshape(simultaneous realization: the operational content of
monotonicity)}]

An admissibility structure \(K\) induces, for each finite \(T \subseteq V\), the registration class
\[
\mathcal{K}(T) \;=\; \{\, \kappa \ \text{sound on the compound view } c_T \;:\; \rho(\kappa) \le \gamma_K(T) \,\},
\]
nonempty and realizing its ceiling by \hyperlink{stmt:14.4}{Lemma~14.4}, and a down-set in resolving power by construction. Say the family \(\{\mathcal{K}(T)\}_T\) is \textbf{closed under restriction} if for all finite \(T' \subseteq T\) and every \(\kappa' \in \mathcal{K}(T')\), the composite \(\kappa' \circ \mathrm{pr}_{T'} \in \mathcal{K}(T)\), where \(\mathrm{pr}_{T'} : \prod_{v \in T} Y_v \to \prod_{v \in T'} Y_v\) is the coordinate projection---\hspace{0pt}i.e., a reading admissible on part of the data remains admissible on all of it. Then
\[
\{\mathcal{K}(T)\}_T \ \text{is closed under restriction} \quad\iff\quad K \ \text{is monotone}.
\]
\end{thmplain}

\begin{proof}

First note \(\kappa' \circ \mathrm{pr}_{T'}\) is sound on \(c_T\) (since \(\mathrm{pr}_{T'} \circ c_T = c_{T'}\) and \(\kappa'\) is sound) with \(\rho(\kappa' \circ \mathrm{pr}_{T'}) = \ker(\kappa' \circ c_{T'}) = \rho(\kappa')\). (\(\Leftarrow\)) \(\rho(\kappa' \circ \mathrm{pr}_{T'}) = \rho(\kappa') \le \gamma_K(T') \le \gamma_K(T)\), the last inequality by monotonicity. (\(\Rightarrow\)) By \hyperlink{stmt:14.4}{Lemma~14.4} choose \(\kappa' \in \mathcal{K}(T')\) with \(\rho(\kappa') = \gamma_K(T')\); closure under restriction gives \(\gamma_K(T') = \rho(\kappa' \circ \mathrm{pr}_{T'}) \le \gamma_K(T)\).
\end{proof}

Monotonicity is thereby not a bare order axiom but a coherence law \emph{across} view sets: it characterizes exactly the ceiling families that a single class of decoders, closed under ignoring part of its input, realizes simultaneously. This---\hspace{0pt}not any inequality on ceilings of unions---\hspace{0pt}is the correct cross-\(T\) constraint at this layer. A subadditivity axiom \(\gamma_K(T \cup T') \le \gamma_K(T) \vee \gamma_K(T')\) would, by contrast, combine with monotonicity to force separability (Definition 15.3) and legislate the registration anomaly of \hyperlink{sec:16}{§16} out of existence; the theory deliberately refrains.

Two canonical instances:

\begin{itemize}
\tightlist
\item
  \textbf{Full admissibility} \(K^{\mathrm{can}}\): \(\gamma_{K^{\mathrm{can}}}(T)=\sigma(T)\). This is the kernel; canonical registration is admissible everywhere. It is monotone.
\item
  \textbf{Ceiling from a resolving structure.} Fix, for each view \(v\), a partition \(C_v\in\mathrm{Part}(D)\) (the finest distinctions the decoder can draw from \(v\)'s presentations) and, for compound views, a ceiling \(C_{\langle T\rangle}\); set \(\gamma_K(T)=C_{\langle T\rangle}\wedge\sigma(T)\). Invariance ceilings (§16) are of this form.
\end{itemize}

\subsecn{3}{14.3}{Separable and joint admissible content}

An admissible interpreter facing a view set \(S\) has two regimes: register each view separately (subject to per-view ceilings) and then combine the pieces, or register the compound view jointly (subject to the compound ceiling). The kernel conflated these; admissibility splits them.

\begin{thmdef}[\hypertarget{stmt:14.5}{Definition~14.5}]

For \(S\subseteq V\) under an admissibility structure \(K\):

\begin{itemize}
\tightlist
\item
  the \textbf{separable content} is
  \[
  \sigma^{\mathrm{sep}}_K(S)=\bigvee_{v\in S}\gamma_K(\{v\});
  \]
\item
  the \textbf{joint content} is
  \[
  \sigma^{\mathrm{jnt}}_K(S)=\gamma_K(S).
  \]
  A question \(q\) is \textbf{separably (resp. jointly) answerable from \(S\) under \(K\)} iff \(\ker(q)\le\sigma^{\mathrm{sep}}_K(S)\) (resp. \(\le\sigma^{\mathrm{jnt}}_K(S)\)).
\end{itemize}
\end{thmdef}

\begin{thmplain}[\hypertarget{stmt:14.6}{Proposition~14.6}\ {\normalfont\upshape(bracketing and ordering)}]

For every \(S\subseteq V\):
\[
\bot\ \le\ \sigma^{\mathrm{sep}}_K(S)\ \le\ \sigma^{\mathrm{jnt}}_K(S)\ \le\ \sigma(S),
\]
where the second inequality requires \(K\) monotone, and the outer two hold for every \(K\). Under full admissibility both middle terms collapse to \(\sigma(S)\).
\end{thmplain}

\begin{proof}

Left and right: (K0)/(K1) give \(\bot\le\gamma_K(\{v\})\) and \(\gamma_K(S)\le\sigma(S)\); joins of terms \(\le\sigma(S)\) stay \(\le\sigma(S)\). Middle: if \(K\) is monotone then \(\gamma_K(\{v\})\le\gamma_K(S)\) for each \(v\in S\), so the join \(\sigma^{\mathrm{sep}}_K(S)\le\gamma_K(S)=\sigma^{\mathrm{jnt}}_K(S)\). Collapse under \(K^{\mathrm{can}}\): \(\gamma(\{v\})=\ker(v)\), whose join is \(\sigma(S)=\gamma(S)\).
\end{proof}

\hyperlink{stmt:14.6}{Proposition~14.6} is the promised bracket: admissible content sits between ``nothing'' and the kernel's canonical content, and the separable/joint pair sits inside that bracket whenever the structure is monotone. The next two sections determine what the kernel theorems become in this bracket, and where they break.

\secn{15}{15}{Fate of the kernel theorems under admissibility}

We take \(\sigma^{\mathrm{sep}}_K\) and \(\sigma^{\mathrm{jnt}}_K\) in turn and ask which of \hyperlink{stmt:6.2}{Theorems~6.2}, 6.4, 6.5, 6.7 survive.

\subsecn{1}{15.1}{The adjunction survives for separable content}

\begin{thmplain}[\hypertarget{stmt:15.1}{Theorem~15.1}\ {\normalfont\upshape(T1 relativized)}]

Fix an admissibility structure \(K\) (neither monotonicity nor the finiteness convention of \hyperlink{sec:14}{§14} is required: separable content depends only on the singleton ceilings, so the statement holds for arbitrary \(S \subseteq V\)). Define \(\tau_K:\mathrm{Part}(D)\to\mathcal P(V)\) by \(\tau_K(\pi)=\{v\in V:\gamma_K(\{v\})\le\pi\}\). Then \((\sigma^{\mathrm{sep}}_K,\tau_K)\) is a Galois connection:
\[
\sigma^{\mathrm{sep}}_K(S)\le\pi\iff S\subseteq\tau_K(\pi).
\]
Consequently the full content of \hyperlink{stmt:6.2}{Theorem~6.2} holds with \(\gamma_K(\{v\})\) in place of \(\ker(v)\): \(\sigma^{\mathrm{sep}}_K\) preserves unions as joins, \(\tau_K\sigma^{\mathrm{sep}}_K\) is a closure operator on \(\mathcal P(V)\), and the admissibly-achievable quotients \(\mathrm{Im}\,\sigma^{\mathrm{sep}}_K\) form a complete lattice.
\end{thmplain}

\begin{proof}

\(\sigma^{\mathrm{sep}}_K(S)=\bigvee_{v\in S}\gamma_K(\{v\})\le\pi\) iff \(\gamma_K(\{v\})\le\pi\) for all \(v\in S\) iff \(S\subseteq\tau_K(\pi)\). The consequences are \hyperlink{stmt:3.4}{Fact~3.4}, exactly as in \hyperlink{stmt:6.2}{Theorem~6.2}, since \(\sigma^{\mathrm{sep}}_K\) is again a join over per-view quantities.
\end{proof}

The reason the adjunction is robust is structural: \(\sigma^{\mathrm{sep}}_K\) is built as a join over singleton contributions \(\gamma_K(\{v\})\), and any such map is automatically a left adjoint. Restricting registration merely replaces each generator \(\ker(v)\) by a coarser generator \(\gamma_K(\{v\})\); the order-theoretic scaffolding is untouched.

Likewise \hyperlink{stmt:6.4}{Theorems~6.4} and 6.5 survive verbatim for \(\sigma^{\mathrm{sep}}_K\): the separably answerable questions form the principal ideal \({\downarrow}\sigma^{\mathrm{sep}}_K(S)\), and a view \(w\) separably helps \(q\) iff \(\gamma_K(\{w\})\) separates some pair unresolved by \(\sigma^{\mathrm{sep}}_K(S)\). The proofs are those of Part I with the generators renamed. The universal property (Theorem 6.7) becomes: \(\sigma^{\mathrm{sep}}_K(S)\) is terminal among consolidations realizable by admissible \emph{per-view} registration.

\subsecn{2}{15.2}{What genuinely changes: joint content is not a join over views}

The joint content \(\sigma^{\mathrm{jnt}}_K(S)=\gamma_K(S)\) is, in general, \textbf{not} expressible as a join of singleton contributions, and this is where the kernel's coherence law fails.

\begin{thmplain}[\hypertarget{stmt:15.2}{Proposition~15.2}\ {\normalfont\upshape(fate of Proposition 6.3)}]

For a monotone \(K\) and \(S\subseteq V\),
\[
\sigma^{\mathrm{sep}}_K(S)\ \le\ \sigma^{\mathrm{jnt}}_K(S),
\]
and the inequality can be strict. Under full admissibility it is an equality (Proposition 6.3). Consequently \(\sigma^{\mathrm{jnt}}_K\) need not preserve unions and need not admit an upper adjoint.
\end{thmplain}

\begin{proof}

The inequality is \hyperlink{stmt:14.6}{Proposition~14.6}. Strictness is exhibited in \hyperlink{sec:16}{§16}. If \(\sigma^{\mathrm{jnt}}_K\) preserved unions it would satisfy \(\gamma_K(S)=\bigvee_{v\in S}\gamma_K(\{v\})=\sigma^{\mathrm{sep}}_K(S)\) for all \(S\); a strict instance refutes this and, by \hyperlink{stmt:3.4}{Fact~3.4}(ii), refutes the existence of an upper adjoint.
\end{proof}

So the kernel identity ``register-then-combine \(=\) combine-then-register'' (Proposition 6.3), which underwrote both the adjunction for joint content and the coincidence of the separable and joint notions, is exactly the statement that \(K\) is \textbf{separable}:

\begin{thmdef}[\hypertarget{stmt:15.3}{Definition~15.3}]

\(K\) is \textbf{separable at \(S\)} if \(\sigma^{\mathrm{sep}}_K(S)=\sigma^{\mathrm{jnt}}_K(S)\), and \textbf{separable} if this holds for all \(S\). The \textbf{registration anomaly} at \(S\) is, for monotone \(K\), the interval
\[
\Delta_K(S)=\big[\,\sigma^{\mathrm{sep}}_K(S),\ \sigma^{\mathrm{jnt}}_K(S)\,\big]\subseteq\mathrm{Part}(D),
\]
nontrivial (non-degenerate) exactly when the endpoints disagree. The interval typing is deliberate: intervals in a complete lattice restrict, intersect, and compose, which is the raw material an obstruction theory needs (§16.4), whereas a bare pair has no algebra. For non-monotone \(K\) the two contents may be incomparable, and the anomaly is then recorded as the unordered pair of endpoints.
\end{thmdef}

\begin{thmplain}[\hypertarget{stmt:15.4}{Corollary~15.4}\ {\normalfont\upshape(finite-view adjunction criterion)}]

Let \(V\) be finite and \(K\) monotone. The following are equivalent: (i) \(K\) is separable; (ii) \(\sigma_K^{\mathrm{jnt}}:\mathcal P(V)\to\mathrm{Part}(D)\) preserves all unions, including the empty union; (iii) it has a right adjoint; (iv) it equals \(\sigma_K^{\mathrm{sep}}\).
\end{thmplain}

\begin{proof}

Separability is (iv). The singleton-generated formula in (iv) preserves unions and has the right adjoint of \hyperlink{stmt:15.1}{Theorem~15.1}. Union preservation gives that formula by expressing any subset as the union of its singletons. A left adjoint preserves unions. Grounding follows from the raw-content bound at the empty set. For infinite \(V\), the finite-subset version characterizes finite-union preservation; the right-adjoint statement applies to the arbitrary-join extension on \(\mathcal P(V)\), not in general to a map with domain \(\mathcal F(V)\).
\end{proof}

The fate of T1 is thus sharp. \textbf{For separable content the adjunction always holds; for joint content it holds iff the admissibility structure is separable, and the obstruction is exactly the registration anomaly \(\Delta_K\).} Full admissibility (the kernel) is separable, which is why Part I never saw the anomaly. Any strictly sub-canonical \(K\) that reads the joint more finely than the combination of its parts breaks separability---\hspace{0pt}and that, far from being pathological, is the generic and interesting case, as the next section shows.

\secn{16}{16}{The registration anomaly, and where it points}

\subsecn{1}{16.1}{A worked non-separable structure: invariance and parity}

We give an explicit \(S\) and \(K\) with \(\sigma^{\mathrm{sep}}_K(S)=\bot\) but \(\sigma^{\mathrm{jnt}}_K(S)\) nontrivial---\hspace{0pt}a maximal anomaly, in which separate admissible registration recovers \emph{nothing} while joint admissible registration answers a real question.

\textbf{Setup.} Let \(D=\{00,01,10,11\}\), thought of as pairs of bits, and let the group \(G=\{\mathrm{id},\phi\}\) act on \(D\) by the global bit-flip \(\phi:00\leftrightarrow11,\ 01\leftrightarrow10\). Two views expose the coordinates,
\[
v(b_1b_2)=b_1,\qquad w(b_1b_2)=b_2,
\]
with \(\ker(v)=\{00\,01\mid10\,11\}\) and \(\ker(w)=\{00\,10\mid01\,11\}\), so \(\sigma(\{v,w\})=\top\). The target question is parity, \(q(b_1b_2)=b_1\oplus b_2\), whose content is \(\ker(q)=\{00\,11\mid01\,10\}\).

\textbf{Admissibility by invariance.} Let the decoder be required to produce readings invariant under \(G\): an admissible registered partition must have \(G\)-stable blocks. The finest \(G\)-invariant partition of \(D\) is the orbit partition
\[
P_G=\{\,00\,11\mid01\,10\,\},
\]
and we take the invariance ceiling \(\gamma_K(T)=P_G\wedge\sigma(T)\) (an instance of the resolving-structure ceiling of \hyperlink{sec:14.2}{§14.2}; one checks (K0)--(K2) and monotonicity hold, since \(\wedge\sigma(T)\) is monotone in \(T\) and \(P_G\) is fixed).

\textbf{Separate content vanishes.} For the single view \(v\),
\[
\gamma_K(\{v\})=P_G\wedge\ker(v)=\{00\,11\mid01\,10\}\wedge\{00\,01\mid10\,11\}.
\]
The meet is the coarsest common coarsening: \(00\sim11\) (from \(P_G\)) and \(00\sim01\), \(10\sim11\) (from \(\ker v\)) chain all four elements together, so \(\gamma_K(\{v\})=\bot\). By symmetry \(\gamma_K(\{w\})=\bot\), hence
\[
\sigma^{\mathrm{sep}}_K(\{v,w\})=\bot\vee\bot=\bot.
\]
Neither coordinate, read invariantly, distinguishes anything: a flip-invariant reading of \(b_1\) alone cannot tell \(b_1=0\) from \(b_1=1\), since \(\phi\) swaps them.

\textbf{Joint content answers parity.} For the compound view \(\langle v,w\rangle\), \(\sigma(\{v,w\})=\top\), so
\[
\sigma^{\mathrm{jnt}}_K(\{v,w\})=\gamma_K(\{v,w\})=P_G\wedge\top=P_G=\{00\,11\mid01\,10\}=\ker(q).
\]
Parity is jointly answerable under \(K\), though nothing is separably answerable. The anomaly is maximal:
\[
\Delta_K(\{v,w\})=[\,\bot,\ \ker(q)\,],\qquad \sigma^{\mathrm{sep}}_K(\{v,w\}) < \sigma^{\mathrm{jnt}}_K(\{v,w\}).
\]

The phenomenon is the informational core of parity secret-sharing \autocite{shamir1979}: each share (coordinate) individually carries nothing about the secret (parity), while the shares jointly determine it. The kernel could not express this, because canonical registration is separable; admissibility is precisely the added structure that makes ``jointly informative, separately null'' a representable---\hspace{0pt}and quantifiable---\hspace{0pt}state of affairs. In measurement-theoretic terms the invariance ceiling is the \emph{meaningfulness} criterion \autocite{krantz1971,narens2002}: admissible content is invariant content, and the anomaly states that meaningful joint content can strictly exceed the join of meaningful marginal contents.

\subsecn{2}{16.2}{The opposite regime, and why monotonicity is the divide}

Strictness in \hyperlink{stmt:15.2}{Proposition~15.2} assumed monotonicity, which delivered \(\sigma^{\mathrm{sep}}_K\le\sigma^{\mathrm{jnt}}_K\). Non-monotone admissibility can invert it. Consider a decoder with an \textbf{output budget}: it emits at most one bit regardless of input, so every admissible reading has at most two blocks. Note that the budget constraint by itself carves out a down-set of \(\mathrm{Part}(D)\) that is \emph{not} a principal ideal---\hspace{0pt}there is no finest two-block partition below \(\top\)---\hspace{0pt}so a budget-limited decoder is modelled, within \hyperlink{stmt:14.3}{Definition~14.3}, by an admissibility structure that makes a definite choice of ceiling inside the budget. (The awkwardness is typal rather than substantive: down-sets of \(\mathrm{Part}(D)\) form a complete lattice under inclusion, so budget classes are first-class content objects in the ideal completion, with ceilings---\hspace{0pt}principal ideals---\hspace{0pt}as the representable case; the graded theory meets the same issue twice more, and the consolidation is made in Part III, \hyperlink{sec:24}{§24}.) In the parity setup, take \(\gamma_K(\{v\})=\ker(v)\) and \(\gamma_K(\{w\})=\ker(w)\) (each already two-block), and for the compound view let the decoder spend its bit on the first coordinate: \(\gamma_K(\{v,w\})=\ker(v)\). Axioms (K0)--(K2) hold, but monotonicity fails: \(\gamma_K(\{w\})=\ker(w)\not\le\ker(v)=\gamma_K(\{v,w\})\). Then
\[
\sigma^{\mathrm{sep}}_K(\{v,w\})=\ker(v)\vee\ker(w)=\top,
\qquad
\sigma^{\mathrm{jnt}}_K(\{v,w\})=\ker(v)<\top:
\]
separate content strictly exceeds joint content.

The two regimes sit on opposite sides of the kernel's equality:
\[\begin{gathered}\underbrace{\sigma^{\mathrm{sep}}_K\ \le\ \sigma^{\mathrm{jnt}}_K}_{\text{monotone: guaranteed (Prop. 14.6)}} \\ \text{vs.} \\ \underbrace{\sigma^{\mathrm{jnt}}_K\ \le\ \sigma^{\mathrm{sep}}_K}_{\text{non-monotone: possible, not forced}},\end{gathered}\]
with equality---\hspace{0pt}Proposition 6.3---\hspace{0pt}exactly at separable structures, of which full admissibility is one. The precise statement is asymmetric: \textbf{monotonicity rules out the resource-limited inversion} \(\sigma^{\mathrm{jnt}}_K < \sigma^{\mathrm{sep}}_K\) (Proposition 14.6); when monotonicity fails, budget-type inversions become possible, though not forced---\hspace{0pt}a non-monotone structure may order the two contents either way, leave them equal, or leave them incomparable, and may do so differently at different view sets. Both regimes are real; the kernel occupies the seam between them.

\subsecn{3}{16.3}{Un-collapsing encoded / extractable / used}

The anomaly lets us separate notions the kernel fused. Reading Part I's canonical content as \emph{what is present in principle} and admissible content as \emph{what a legitimate decoder recovers}, \hyperlink{stmt:14.6}{Proposition~14.6} stratifies content over \(S\):
\[
\underbrace{\sigma(S)}_{\text{encoded}}\ \ge\ \underbrace{\sigma^{\mathrm{jnt}}_K(S)}_{\text{extractable}}\ \ge\ \underbrace{\sigma^{\mathrm{sep}}_K(S)}_{\text{separably extractable}}\ \ge\ \underbrace{\vphantom{\sigma^{\mathrm{jnt}}}\ \cdots\ }_{\text{used}}
\]
(the monotone case). \emph{Encoded} is what the views draw at all; \emph{extractable} is what admissible joint registration recovers; \emph{separably extractable} is what admissible per-view registration recovers; \emph{used}---\hspace{0pt}content that additionally drives a downstream process---\hspace{0pt}is a further coarsening requiring a model of that process and is deferred (it modifies neither P1--P6 nor the admissibility layer but adds dynamics; \hyperlink{sec:8}{§8}, beyond H2). The distinctions \emph{encoded \(\ne\) extractable \(\ne\) separably extractable} are now theorems, witnessed by strict instances such as \hyperlink{sec:16.1}{§16.1}, rather than informal slogans. This is the stratification the admissibility layer was introduced to provide: in particular it separates ``the representation contains the distinction'' (\(\le\sigma\)) from ``an admissible decoder can draw it'' (\(\le\sigma^{\mathrm{jnt}}_K\)), which is the exact ambiguity that undermines interpretability claims resting on unconstrained probes.

\subsecn{4}{16.4}{The anomaly as the entry point to gluing}

For monotone \(K\), \(T\mapsto\gamma_K(T)\) is covariant in the view set. Separability says that its singleton contributions jointly attain its value on \(T\). Locally admissible readings always combine under monotonicity; the possible failure is that those combinations do not dominate all joint readings (Proposition 36.3). Thus the interval \(\Delta_K\) is a cofinality defect, not by itself a failure of the sheaf axiom for a Set-valued presheaf.

Part V develops positive union covers, orbit graphs, and witness sets. Appendix D.8 states the additional distributivity condition needed for join covers on achievable partitions to form a site. Appendix D.12 shows why the parity admissibility anomaly does not imply contextuality of the ordinary marginal empirical model. These distinctions preserve the local-to-global motivation without identifying different gluing problems.

\secn{17}{17}{Established results and scope}

The exact admissibility theory proves realizability, projection coherence, bracketing, and the separability criterion. The anomaly is a failure of local readings to be cofinal among joint readings, not a failure to construct their combination. Its graph and positive-cover analysis is developed in Part V; Appendix D.8 identifies exactly when join covers on achievable partitions form a site.
